\section{논의 및 결론}
\label{sec:discussion}

\subsection{시공간 데이터베이스 표현의 의미}

이 논문의 출발점은 PUBG 스쿼드 매치를 하나의 최종 순위 기록으로 압축하지 않고,
팀이 페이즈마다 위험집합에 남거나 탈락하는 과정으로 읽는 것이었다. 이 관점에서
PUBG 텔레메트리는 단순한 게임 로그가 아니라 이동 객체, 시간변화 영역, 사건,
위험집합으로 구성된 시공간 데이터베이스가 된다. 이 재구성은 예측 문제의 정의를
바꾼다. 최종 순위가 이미 결정된 뒤 집계값으로 맞히는 문제가 아니라, 특정 페이즈
종료 시점까지 관측된 정보만으로 다음 페이즈의 탈락 위험을 예측하는 문제가 된다.

여기서 시공간 데이터베이스 표현은 입력 변수를 많이 만드는 절차가 아니다. 어떤
정보가 어느 시점까지 관측 가능하고, 어떤 팀이 아직 위험집합에 남아 있으며, 어떤
공간 영역과 사건이 그 시점의 위험을 구성하는지 정하는 절단 규칙이다. 이 절단
규칙이 있어야 PUBG 텔레메트리를 다음 페이즈 생존 예측 문제로 바꿀 수 있다.

\subsection{초기 실험이 보여준 우선순위}

초기 비교 실험에서 가장 안정적으로 작동한 정보는 팀 생존 여력이었다. 평균 체력,
최소 체력, 생존 인원, 차량 이용 여부는 모든 맵/플레이 모드 조건에서 가장 낮은
손실을 보였다. 이는 전투 생존 게임에서 자연스러운 결과다. 남은 체력과 생존 인원은
다음 페이즈를 버틸 수 있는지를 직접 나타내고, 차량 이용 여부도 이동 가능한 자원을
가리킨다. 그래서 이 결과를 새로운 발견처럼 과장해서는 안 된다. 이 실험에서 확인한
것은 직접 생존 상태 변수가 강하다는 점이고, 논문의 시공간 표현 질문은 그 다음에
온다.

그 다음 질문은 직접 생존 상태 없이도 시공간 이력이 정보를 갖는가 하는 점이다. 이동 이력
조건은 팀 중심 좌표와 속도만 입력으로 사용했는데도 최신 상태 기준모델보다 낮은
손실을 보였다. 팀을 마지막 좌표 하나로 줄이면 위험한 공간에 어떻게 들어갔는지
사라진다. 이 결과는 시공간 이력 보존이 적어도 이 수준에서는 생존 예측에 도움이
되었음을 보여준다. 다만 그 결론을 모든 구성요소로 넓히기는 어렵다.

\subsection{안전구역 표현의 한계}

안전구역은 배틀로얄 장르의 중심적인 구조다. 그런데 이번 실험은 안전구역 맥락이
단독 표현으로 강한 예측 신호를 낸다는 주장을 지지하지 않았다. 더 불편한 부분은
이동 이력에 안전구역 변수를 단순 결합한 조건도 이동 이력만 사용한 조건보다 좋아지지
않았다는 점이다. 이 결과를 피하지 않는 편이 낫다. 현재 방식은 안전구역을 잘 쓴
모델이라기보다, 안전구역 변수를 좌표 옆에 붙여 본 수준에 가깝다. 후속 모델에서는
안전구역을 독립 입력 블록으로 두기보다 이동 이력과 팀 생존 여력의 의미를 바꾸는
맥락으로 다루어야 한다. 같은 속도라도 남은 페이즈 시간, 안전구역 반경, 주변 팀
밀도에 따라 다른 탈락 위험을 가질 수 있다.

\subsection{한계}

현재 결과는 초기 비교 실험에 기반한다. 또한 본 실험은 공개 telemetry 전체 feature를
사용한 official telemetry-derived full baseline과 STDB 전체 표현을 비교한 것이 아니다.
ST 관점에서 추출한 13개 snapshot feature 내부에서 어떤 block이 먼저 작동하는지
확인한 ablation으로 해석해야 한다. 각 조건은 같은 데이터 분할에서 비교되었지만
seed 반복, 긴 학습, 하이퍼파라미터 조정은 아직 충분하지 않다. 또한 표현 구성요소
비교에는 제한 표현, 단순 결합 표현, 압축 기준모델이 함께 포함되어 있다. 이 때문에
이를 완전한 단계적 추가 실험으로 해석해서는 안 된다. 전체 팀 상태 스냅샷 조건은
입력 차원이 더 크기 때문에 짧은 학습에서 불리했을 가능성이 있다. 또한 본 버전의
정량 결과는 다음 페이즈 탈락 위험에 한정된다. 최종 순위와
top-\(k\) 진입은 같은 생존 과정의 장기 결과지만, 여기서는 직접 평가하지 않았다.
마지막으로 PUBG 텔레메트리는 지형, 엄폐, 시야, 플레이어 숙련도, 팀 의사소통 같은
중요한 조건을 직접 관측하지 못한다.

\subsection{결론}

PUBG 공식 텔레메트리는 최종 순위 행 데이터로만 남길 필요가 없다. 이동 객체,
시간변화 영역, 시공간 사건, 위험집합으로 재구성하면 페이즈별 탈락
위험이라는 예측 문제가 만들어진다. 초기 실험에서 가장 강한 신호는 팀 생존 여력이었지만,
이는 전투 생존 게임에서 예상 가능한 직접 상태 변수의 효과다. 시공간 표현 관점에서
더 중요한 결과는 이동 이력 조건이 좌표와 속도만으로도 최신 상태 기준보다 낮은
손실을 보였다는 점이다. 안전구역은 단독 또는 단순 결합 조건에서 충분히 활용되지
않았다. 다음 단계에서는 안전구역을 이동 이력과 팀 생존 상태를 조건화하는 상호작용
기반 시공간 표현으로 확장해야 한다.
